% Source: https://artofproblemsolving.com/community/c3h2510838_what_up_with_mass_points
% Stable identity: aops-topic-2510838-q01
\textbf{Why ordinary mass points get stuck.} The transversal $EF$ cuts two sides through $A$, so $A$ would need one mass to balance $B$ about $F$ (where $AF:FB=1:3$ needs mass $3$ at $A$ for mass $1$ at $B$) and another to balance $C$ about $E$ (where $AE:EC=1:1$ needs equal masses). The fix is to \emph{split} the mass at $A$.

\textbf{Split masses.} Put mass $1$ at $B$ and $1$ at $C$, so $D$, the midpoint of $BC$, gets mass $2$.
\begin{itemize}
\item On $AB$: $AF:FB=1:3$, so the part of $A$'s mass used with $B$ is $3$, and $F$ gets $3+1=4$.
\item On $AC$: $E$ is the midpoint, so the part of $A$'s mass used with $C$ is $1$, and $E$ gets $1+1=2$.
\end{itemize}
So $A$ has total mass $3+1=4$.

\begin{asy}
usepackage("amsmath");
size(5.5cm);
defaultpen(fontsize(10pt));
pair B=(0,0),C=(6,0),A=(1.6,4.2);
pair D=(B+C)/2, E1=(A+C)/2, F=A+(B-A)/4;
pair X=extension(E1,F,A,D);
draw(A--B--C--cycle,linewidth(0.9));
draw(A--D,blue); draw(E1--F,red);
dot(A);dot(B);dot(C);dot(D);dot(E1);dot(F);dot(X);
label("$A$",A,N);label("$B$",B,SW);label("$C$",C,SE);label("$D$",D,S);label("$E$",E1,NE);label("$F$",F,W);label("$X$",X,dir(-20));
label("$3+1$",A+(0.7,0.1),fontsize(8pt)+deepgreen); label("$1$",B+(-0.3,0.35),fontsize(8pt)+deepgreen); label("$1$",C+(0.3,0.35),fontsize(8pt)+deepgreen);
label("$2$",D+(0.35,0.3),fontsize(8pt)+deepgreen); label("$2$",E1+(0.1,-0.45),fontsize(8pt)+deepgreen); label("$4$",F+(-0.2,0.5),fontsize(8pt)+deepgreen);
\end{asy}

The centre of mass of the whole system lies on $AD$ (mass $4$ at $A$, mass $2$ at $D$) and on $EF$ (mass $2$ at $E$, mass $4$ at $F$), so it is $X$. On $AD$ the distances are inversely proportional to the masses:
\[AX:XD=2:4=1:2,\qquad\dfrac{AX}{AD}=\boxed{\dfrac13}.\]
(Also $EX:XF=4:2=2:1$.)

\textbf{Check with vectors.} Take $A$ as origin and write $\vec b=\overrightarrow{AB}$, $\vec c=\overrightarrow{AC}$. Then $D=\dfrac{\vec b+\vec c}{2}$, $E=\dfrac{\vec c}2$, $F=\dfrac{\vec b}4$. A point of $AD$ is $\lambda\dfrac{\vec b+\vec c}{2}$ and a point of $EF$ is $(1-\mu)\dfrac{\vec c}2+\mu\dfrac{\vec b}4$. Comparing coefficients of $\vec b$ and $\vec c$: $\dfrac\lambda2=\dfrac\mu4$ and $\dfrac\lambda2=\dfrac{1-\mu}2$, so $\mu=\dfrac23$ and $\lambda=\dfrac13$, confirming $\dfrac{AX}{AD}=\dfrac13$.
